%! AP Statistics — Unit 3, Lesson 3.7 — Warm-Up KEY
\documentclass[10pt]{article}
\usepackage{apstats-article}
\usepackage{apstats-key}

\begin{document}

\pageheader{Unit 3, Lesson 3.7}{Warm-Up --- Answer Key}
\begin{tabularx}{\linewidth}{X X X}
Name:\;\ans{Key} & Date:\; & Period:\;
\end{tabularx}

\vspace{0.3in}

{\large\bfseries\color{navy} Part 1 \quad Spiral Review}

\vspace{0.12in}

\begin{enumerate}

    \item
    \renewcommand{\arraystretch}{1.8}
    \begin{tabularx}{\linewidth}{@{} X p{2.8cm} p{3.4cm} @{}}
    \textbf{Scenario} & \textbf{Design} & \textbf{Brief justification} \\
    \hline
    Test new fertilizer on 30 identical corn plants. &
    \ans{CRD} & \ans{Units are homogeneous; no blocking needed.} \\
    Test two teaching methods on 40 students varying in reading ability. &
    \ans{Block design} & \ans{Block on reading ability to reduce variability.} \\
    \end{tabularx}

    \vspace{0.18in}

    \item
    \begin{enumerate}[label=(\alph*), itemsep=4pt]
        \item \ans{Study B --- random assignment was used to assign drug or placebo.}
        \item \ans{Study B --- the sample was randomly selected from the patient
        population, so results generalize to that population.}
        \item \ans{Study B used both random assignment (allowing causal conclusions)
        and random selection (allowing generalization). Study A used neither, so
        results can only describe those specific 50 volunteers.}
    \end{enumerate}

\end{enumerate}

\vspace{0.2in}

{\large\bfseries\color{navy} Part 2 \quad Looking Ahead}

\vspace{0.14in}

\noindent\textcolor{slate}{\small Student responses will vary. Accept: random assignment
controls for confounding variables by distributing them equally across groups; it does
not tell us about the broader population, which requires random selection.}

\end{document}
